\section{Finite actual geometric representatives of every intrinsic ample class}
\label{sec:native-numerical-positive-ample-sums}

\begin{theorem}[Positive finite original ample-line representatives {\cite[\S1.2]{BJ22}}]
\label{thm:native-numerical-positive-ample-sums}
\lean{EnvelopeProperty.NativeRealNumericalClasses.exists_positive_ample_line_sum, EnvelopeProperty.NativeRealNumericalClasses.curveDegree_ampleCone_pos}
\leanok
\uses{def:native-numerical-ample-cone,thm:native-numerical-ample-classes}
Every class in the original intrinsic ample cone has an internally
produced nonempty finite family of actual geometrically ample original
lines $A_i$ and strictly positive real coefficients $a_i$, with
\[
 \theta=\sum_i a_i[A_i]
\]
in the same original numerical quotient. Consequently every such class
has strictly positive genuine full Euler degree on every original
closed integral curve.
\end{theorem}
\begin{proof}
The set of nonempty finite positive sums of actual ample-line classes
is a convex cone. Positive scaling multiplies every coefficient by the
same positive real scalar. Addition concatenates the two finite
families, preserving their actual lines, their geometric ampleness and
strictly positive coefficients. Each generating original ample-line
class is represented by its one-element family with coefficient one.
The minimality of the original convex-cone hull therefore yields the
stated representation for every class in that same hull. No finite
dimensionality or selected numerical presentation is assumed.

For strict curve positivity, the classes having positive genuine degree
on every original closed curve form a convex cone: the original degree
map is real linear, a positive scalar preserves strict positivity, and
the sum of two strictly positive degrees is strictly positive. Actual
geometric ample-line positivity puts every original generator in this
cone. Hull minimality proves the strict positivity of the entire
original ample cone, with all original full Euler curve tests retained.
\end{proof}
